% ===================================================================
%  Human Brain Mapping — Technical Report
%  Scanner Vendor Harmonisation of Spinal Cord qMRI Biomarkers
% ===================================================================
\documentclass[12pt,a4paper]{article}

% Packages
\usepackage[utf8]{inputenc}
\usepackage[T1]{fontenc}
\usepackage{amsmath,amssymb}
\usepackage{graphicx}
\usepackage{booktabs}
\usepackage[margin=2.5cm]{geometry}
\usepackage{hyperref}
\usepackage{setspace}
\usepackage{caption}
\usepackage{enumitem}
\usepackage{float}
\usepackage{lineno}
\usepackage{multirow}
\usepackage{siunitx}
\usepackage{tabularx}
\usepackage{xurl}
\usepackage{pdflscape}
\setlength{\emergencystretch}{3em}

\hypersetup{colorlinks=true,linkcolor=blue,citecolor=blue,urlcolor=blue}
\onehalfspacing

\begin{document}
\linenumbers

% ===================================================================
\section*{Abstract}

\noindent\textbf{Objectives:} Quantitative spinal cord MRI biomarkers exhibit systematic vendor-dependent variability that limits multi-centre pooling, yet no systematic harmonisation benchmark exists for spinal cord data.

\noindent\textbf{Methods:} Eight qMRI biomarkers (T2w-CSA, GM-CSA, MTR, MTsat, FA, MD, AD, RD) were extracted at C2--C3 from the spine-generic dataset (267 subjects, 43 sites, 3 vendors; site and vendor perfectly nested). Five methods---ComBat, ComBat-joint, CovBat, RELIEF, LME/FE---were evaluated on four endpoints in healthy-control ($N=188$) and pathology-inclusive ($N=246$) cohorts: vendor-effect removal, biological-signal recovery, subject-level preservation, and machine-learnability of the vendor signature. Leave-one-vendor-out validation tested generalisability.

\noindent\textbf{Results:} Baseline vendor effects were substantial ($\eta^2=0.31$; PERMANOVA $R^2=0.315$, $p=1\times10^{-4}$, the permutation lower bound with 9,999 permutations). All methods reduced vendor $\eta^2$ to ${\sim}10^{-4}$ (per-metric reductions of 92.1--100.0\%; post-harmonisation values lie far below the $\eta^2\approx0.011$ expected by chance)---a near-degenerate criterion, since vendor is the harmonisation batch variable. In leakage-free fold-internal evaluation, random-forest vendor detection fell from AUC~0.99 (unharmonised) to 0.63--0.81 across methods in the pathology-inclusive cohort (HC: 0.67--0.85; permutation reference 0.50), demonstrating substantial but incomplete removal of the multivariate vendor signature, lowest for CovBat and highest for LME/FE; pathology classification concurrently improved (AUC 0.52 to 0.62, ComBat; paired $p<0.001$). Age associations absent at baseline (all $p\geq0.19$) became significant in diffusion metrics after harmonisation (MD/AD, $p\leq0.0019$ in the pathology-inclusive cohort). Methods diverged on subject preservation (LME/FE: $r=0.82$; CovBat: $r=0.78$), forming a Pareto frontier; within-resample paired bootstrap confirmed the preservation hierarchy (8/10 and 9/10 method pairs significantly separated in HC and ALL, respectively). LOVO validation showed ComBat/ComBat-joint generalised without significant degradation (PERMANOVA $p\geq0.97$ in all folds), while RELIEF and LME/FE failed ($\eta^2\approx0.2$, $p<0.01$). Three methodological constraints---vendor-site nesting, LME fixed-effects fallback (56\% of internal cases), and RELIEF ICA inactivation---converge on a single structural limitation: only 3 vendors in a perfectly nested 43-site design.

\noindent\textbf{Conclusion:} Harmonisation is recommended for pooled multi-vendor spinal cord qMRI. ComBat-joint, which harmonises all eight biomarkers jointly, provided the most stable generalisation to unseen scanner vendors (LOVO/internal $\eta^2$ ratio 0.9 vs.\ 6.0 for per-metric ComBat) together with the highest subject-level preservation among ComBat-family methods, and is preferred for broad deployment. Per-metric ComBat remains appropriate when individual biomarkers are harmonised in isolation, and CovBat/RELIEF suit pooled-group analyses within fixed vendor sets.

\vspace{4mm}
\noindent\textbf{Key Points:}
\begin{enumerate}[nosep,leftmargin=*]
  \item Scanner vendor effects dominate biological variation in spinal cord quantitative MRI by orders of magnitude, making harmonisation essential for reliable multi-centre biomarker comparison.
  \item ComBat-family harmonisation removes most machine-detectable vendor information (leakage-free AUC 0.99 to ${\approx}0.70$ for ComBat and ComBat-joint; chance 0.50) while unmasking age- and pathology-related biological signal previously obscured by vendor variance.
  \item Leave-one-vendor-out validation reveals that latent-structure (RELIEF) and random-effects (mixed-model) approaches fail to generalise to new scanner vendors; among the methods that did transfer, joint ComBat estimation was the most stable (LOVO/internal $\eta^2$ ratio 0.9 vs.\ 6.0 for per-metric ComBat), supporting ComBat-joint for multi-centre consortia.
\end{enumerate}

\vspace{4mm}
\noindent\textbf{Keywords:} Spinal cord; Quantitative MRI; Harmonisation; ComBat; Multicentre study; Reproducibility; Scanner vendor effects

\vspace{4mm}
\noindent\textbf{Abbreviations:} AD = axial diffusivity; ComBat = Combatting batch effects; CovBat = Covariance Batch harmonisation; CSA = cross-sectional area; DCM = degenerative cervical myelopathy; FA = fractional anisotropy; HC = healthy control; ICA = independent component analysis; LME/FE = linear mixed-effects / fixed-effects hybrid; LOVO = leave-one-vendor-out; MD = mean diffusivity; MT = magnetisation transfer; MTR = magnetisation transfer ratio; MTsat = magnetisation transfer saturation; PERMANOVA = permutational multivariate analysis of variance; qMRI = quantitative magnetic resonance imaging; RD = radial diffusivity; RELIEF = REmoval of Latent Inter-scanner Effects through Factorization

% ===================================================================
\section{Introduction}

Quantitative MRI (qMRI) biomarkers of the spinal cord---cross-sectional area (CSA), magnetisation transfer ratio (MTR), magnetisation transfer saturation (MTsat), and diffusion tensor metrics (FA, MD, AD, RD)---provide non-invasive windows into atrophy, demyelination, and axonal integrity in multiple sclerosis, degenerative cervical myelopathy, and spinal cord injury [1--3]. Multi-centre consortia are increasingly adopting these biomarkers, but systematic vendor-dependent variability limits cross-site pooling [4,5]. A recent pan-European survey by the European Society of Radiology documented marked heterogeneity in MRI quality assurance and control practices across 37 countries [6], underscoring the urgency of robust harmonisation.

Even under standardised acquisition protocols, differences in gradient hardware, RF coil geometry, and vendor-specific reconstruction pipelines introduce offsets that often exceed the biological effect sizes of interest [4,7]. Statistical post-hoc harmonisation methods---ComBat [8--10], CovBat [11], RELIEF [12], and linear mixed-effects models (LME) [13]---have been extensively benchmarked in brain imaging and reviewed in the broader context of multicentre radiomics [14] but remain uncharacterised for spinal cord qMRI. Fortin et al.\ first applied ComBat to multi-site DTI [10], establishing a framework later extended to cortical thickness [9], longitudinal harmonisation [15], covariance harmonisation [11], latent-factor removal [12], and out-of-sample harmonisation for new sites [16]. Middleton et al.\ demonstrated ComBat's effectiveness for spinal cord DTI [17], and Jodoin et al.\ provided deployment guidance for diffusion MRI [18], but no study has systematically compared the full ComBat family across a comprehensive spinal cord biomarker panel using unified multi-endpoint evaluation.

The spinal cord presents distinct harmonisation challenges: small per-site sample sizes (1--8 subjects per site in the complete-case sample, median 5) challenge empirical Bayes shrinkage; strong inter-metric correlations (diffusion metrics are intrinsically coupled; MTR and MTsat share acquisition-dependent variance) motivate multivariate evaluation; and the clinical need to preserve disease-related variance creates a tension not previously quantified.

Here we benchmark five methods on eight spinal cord qMRI biomarkers using four orthogonal endpoints: vendor-effect removal, biological-signal recovery, subject-level preservation, and machine-learnability of the vendor signature. We extend prior work through: (i) formal comparison of the full ComBat family (including covariance and latent-factor extensions); (ii) multivariate evaluation via PERMANOVA; (iii) a Pareto-frontier analysis of the decontamination--preservation trade-off; and (iv) leave-one-vendor-out external validation to test generalisability.

% ===================================================================
\section{Materials and Methods}

\subsection{Dataset and Cohorts}

The spine-generic multi-subject dataset (release r20231212) [1] comprises 267 subjects scanned at 43 sites across three vendors (GE Healthcare: 37 subjects at 6 sites; Philips: 50 subjects at 8 sites; Siemens Healthineers: 180 subjects at 29 sites) using a standardised consensus protocol [4]. Site and vendor are perfectly nested (each site operates a single scanner from one vendor); reported vendor effects thus represent the composite ``vendor-in-site'' effect. The complete-case analysis sample comprised 40 sites (GE 5, Philips 8, Siemens 27). Demographics were balanced across vendors in the complete-case analysis sample (Table~1; all $p\geq0.203$).

\begin{table}[H]
\centering
\small
\setlength{\tabcolsep}{4pt}
\caption{Demographic and acquisition characteristics by vendor (HC complete-case analysis sample, $N=188$). Continuous variables: mean $\pm$ SD, one-way ANOVA. Categorical variables: count (\%), $\chi^{2}$ test. Site and vendor are perfectly nested (40 sites across 3 vendors in the analysis sample; 43 sites in the full dataset).}
\label{tab:demographics}
\begin{tabular}{lrrrrr}
\toprule
Characteristic & GE & Philips & Siemens & Total & $p$-value \\
\midrule
$N$ & 25 & 33 & 130 & 188 & --- \\
Sites ($n$) & 5 & 8 & 27 & 40 & --- \\
Subjects per site, median [range] & 5 [2--8] & 4 [2--6] & 5 [2--8] & 5 [2--8] & --- \\
Age, years (mean $\pm$ SD) & 27.0 $\pm$ 4.9 & 27.9 $\pm$ 4.1 & 29.0 $\pm$ 6.0 & 28.5 $\pm$ 5.6 & 0.203 \\
Age range & 19--39 & 23--38 & 20--52 & 19--52 & --- \\
Sex = F & 13 (52.0\%) & 18 (54.5\%) & 60 (46.2\%) & 91 (48.4\%) & 0.640 \\
Sex = M & 12 (48.0\%) & 15 (45.5\%) & 70 (53.8\%) & 97 (51.6\%) & --- \\
Height, cm & 170.1 $\pm$ 8.7 & 171.3 $\pm$ 9.9 & 173.2 $\pm$ 10.1 & 172.5 $\pm$ 9.9 & 0.277 \\
Weight, kg & 66.4 $\pm$ 13.6 & 66.5 $\pm$ 13.6 & 67.9 $\pm$ 13.1 & 67.5 $\pm$ 13.2 & 0.781 \\
Scanner models ($n$) & 3 & 4 & 5 & 12 & --- \\
\bottomrule
\end{tabular}
\end{table}

Two pre-specified cohorts were analysed: HC ($N=188$ complete-case healthy controls) isolates pure vendor-driven variance; ALL ($N=246$ complete-case, including 56 mild cervical compression and 2 DCM) approximates real-world deployment. Pathology distribution was balanced across vendors ($\chi^2\ p=0.661$). Detailed cohort flow and missingness patterns are provided in Supplementary Material~S1.

\subsection{Biomarker Extraction and Processing}

Eight biomarkers spanning four acquisition families were extracted at the C2--C3 level using the Spinal Cord Toolbox (SCT v7.4) [19]: T2w-CSA (T2-weighted structural), GM-CSA (T2\textsuperscript{*}-weighted), MTR and MTsat (magnetisation transfer), and FA, MD, AD, RD (diffusion tensor imaging). Processing followed the spine-generic pipeline documentation. Automated quality-control metrics excluded acquisitions with segmentation or model-fit failures (see Supplementary Material~S2).

\subsection{Harmonisation Methods}

Five methods were evaluated (full mathematical details in Supplementary Material~S3). In all methods, scanner vendor (GE/Philips/Siemens, three levels) served as the batch variable; because sites are perfectly nested within vendors, vendor-level and site-level offsets are statistically confounded, and the composite ``vendor-in-site'' effect is the harmonisation target:

\textbf{ComBat} [8--10]: Per-metric location-and-scale adjustment with empirical Bayes shrinkage on vendor-specific parameters. Each biomarker harmonised independently.

\textbf{ComBat-joint} [9]: All eight biomarkers harmonised jointly (joint empirical Bayes), preserving cross-metric covariance structure.

\textbf{CovBat} [11]: ComBat-joint followed by principal component (PC) correction---residuals projected onto PCs, each PC harmonised independently, then projected back. PCs explaining at least 95\% of the residual variance were retained: $K=5$ of 8 (cumulative variance 95.9\% in the HC cohort, 95.6\% in the ALL cohort; the first three PCs accounted for 85.6\% and 84.9\%, respectively).

\textbf{RELIEF} [12]: ComBat-joint followed by independent component analysis (ICA) of residuals. Our implementation is a residual-ICA variant of RELIEF (REmoval of Latent Inter-scanner Effects through Factorization), rather than the linked matrix factorisation of the original report; this distinction is revisited in the Limitations. Components with significant vendor association (Benjamini-Hochberg FDR-controlled ANOVA, $q<0.05$) are zeroed. In the present dataset, no component reached significance (all $q\geq0.99$), rendering the pipeline functionally equivalent to ComBat-joint with an additional PCA/ICA round-trip.

\textbf{LME/FE}: Vendor as a random intercept with age and sex fixed effects. When the random-effects covariance matrix was singular (expected with only 3 vendor levels---well below the recommended $k\geq5$ for stable REML estimation), the model fell back to ordinary least squares centring (9 of 16 biomarker$\times$cohort combinations, 56\%; see Results). Only MD and AD achieved stable LME fits in both cohorts. We label this hybrid approach ``LME/FE'' to acknowledge the dual mechanism.

\subsection{Evaluation Endpoints}

Four orthogonal endpoints were evaluated per method per cohort (detailed in Supplementary Material~S4):

\textbf{Vendor-effect removal}: Univariate $\eta^2$ (ANOVA), intraclass correlation coefficient (ICC), and multivariate PERMANOVA $R^2$ (Euclidean distance, 9,999 permutations) [20]. All five methods and the unharmonised baseline---including ComBat and LME/FE, which were refitted on the unified complete-case sample---were evaluated on the same subjects, with per-dimension z-scoring of biomarkers before PERMANOVA. Consequently, the multivariate $R^2$ equals the mean of the per-metric $\eta^2$ for every method (a mathematical identity of z-scored Euclidean PERMANOVA), so univariate and multivariate summaries are directly comparable throughout.

\textbf{Biological-signal recovery}: Type-II semi-partial $R^2$ for age and sex from linear models.

\textbf{Subject-level preservation}: Pearson $r$ between pre- and post-harmonisation per-subject biomarker vectors.

\textbf{Machine-learnability of the vendor signature}: Following the site-detection paradigm of Chen et al.\ [11], a random forest (500 trees, minimum leaf size 5, balanced class weights) was trained to classify scanner manufacturer (Siemens/GE/Philips) from the 8-biomarker profile of each subject, using repeated stratified 5-fold cross-validation (20 repeats). To prevent information leakage, each harmonisation model was refitted on the training folds only and applied to the held-out fold within every cross-validation split; because scanner vendor is observable at deployment, vendor labels were available to the harmonisation step on training and test folds alike, while the random forest itself remained blind to vendor labels. The macro one-vs-rest AUC quantifies residual vendor information detectable by machine learning, a stricter criterion than univariate $\eta^2$ because nonlinear combinations of features can encode vendor structure invisible to marginal statistics. Biological-signal preservation was indexed by the same classifier applied to pathology (healthy control vs disease, ALL cohort, $N=246$ complete cases) and sex labels. A shuffled-label permutation reference calibrated the chance level.

\subsection{Leave-One-Vendor-Out External Validation}

For each method, harmonisation models were trained on subjects from two vendors and applied to the held-out third vendor in three rotation folds. Test-vendor transformation used z-score distribution matching (ComBat, ComBat-joint) or projection onto training PCA/ICA space (CovBat, RELIEF). LME/FE assigned offset=0 to unseen vendors. Three of the four endpoints (vendor-effect removal, biological-signal recovery, and subject-level preservation; machine-learnability and the ALL-cohort pathology/sex preservation indices were not re-evaluated) were evaluated after pooling the harmonised training-vendor and held-out test-vendor data within each fold (Supplementary Material~S5). For ComBat, ComBat-joint, and LME/FE, the test-vendor transformation applies a per-biomarker affine (location--scale) map, so the test-vendor pre/post correlation equals 1.0 by construction; for CovBat and RELIEF, projection onto the training PCA/ICA basis mixes biomarker dimensions, and the test-vendor pre/post correlation fell below 1 ($r=0.85$--$0.97$), indicating genuine distortion of held-out-vendor subject rankings. LOVO preservation statistics for the first group therefore reflect the pooled sample rather than the unseen vendor in isolation. Similarly, the LOVO/internal $\eta^2$ ratio below 1 for ComBat-joint reflects the pooled evaluation sample rather than genuine improvement on unseen vendors.

\subsection{Supplementary Analyses}

Design C balanced-subsample analysis restricted Siemens to 40 randomly drawn subjects per seed (HC cohort, retaining all GE and Philips subjects; $n\approx98$ per subset; 20 random seeds), with all methods refitted per subset using the same canonical implementations as the main analysis, to isolate the effect of vendor imbalance from method performance. Full results in Supplementary Material~S6.

\subsection{Uncertainty Quantification}

Ninety-five per cent confidence intervals for key point estimates were obtained using a stratified nonparametric bootstrap ($B=2{,}000$), resampling subjects with replacement within each scanner-vendor stratum to preserve group proportions. Percentile intervals are reported for $r_{\text{pre,post}}$ and for the baseline vendor effect size. Because marginal CIs for $r_{\text{pre,post}}$ overlap across methods, method pairs were additionally compared with a within-resample paired bootstrap: for each bootstrap sample, all methods' $r$ values were recomputed and their differences tested against zero, which removes between-resample variance and directly targets the pairwise contrast. Bootstrap intervals are not reported for post-harmonisation $\eta^2$ or PERMANOVA $R^2$ because harmonisation methods optimise on the same data, rendering subject-level resampling intervals unreliable for these near-degenerate statistics.

\subsection{Reporting Guidelines and Missing Data}

This study follows STROBE reporting guidance for cross-sectional studies where applicable. The spine-generic dataset is a curated benchmark with pre-processed, quality-controlled biomarker measurements. Subjects with missing biomarker data were excluded via complete-case analysis (HC: $N=188/203$, 7.4\%; ALL: $N=246/267$, 7.9\%); missingness was attributable to acquisition-level exclusions (motion artefacts, segmentation failures, and incomplete sequences) as detailed in Supplementary Material~S1. Sensitivity to missing-data assumptions was assessed through the ALL-cohort analysis, which preserves a broader subject spectrum.

% ===================================================================
\section{Results}

\subsection{Baseline Vendor Effects}

At baseline, vendor explained 31\% of univariate biomarker variance (mean $\eta^2=0.315$, 95\% CI [0.273, 0.381], stratified bootstrap; range 0.05--0.70 across metrics) and 31\% of multivariate structure (PERMANOVA $R^2=0.315$, $p=1\times10^{-4}$, the permutation lower bound with 9,999 permutations; HC cohort). These effects exceeded age-related variance by approximately two orders of magnitude and sex-related variance by one order. AD showed the largest vendor effect ($\eta^2_{\text{HC}}=0.70$), followed by MD ($0.46$) and MTsat ($0.34$); among structural metrics, GM-CSA showed the strongest vendor discrepancy ($\eta^2=0.18$, Philips scanners systematically lower; Supplementary Fig.~S3). In the ALL cohort, baseline vendor effects were comparable or larger ($\eta^2=0.336$; PERMANOVA $R^2=0.336$).

\begin{figure}[H]
\centering
\includegraphics[width=\textwidth]{figures/Fig1.pdf}
\caption{Overview of the study dataset and biomarker panel. \textbf{(a)} Study cohort: the spine-generic multi-subject dataset (release r20231212) [1] comprising 267 subjects scanned at 43 sites on 13 scanner models from three vendors (GE: 37 subjects at 6 sites; Philips: 50 at 8; Siemens: 180 at 29), with the two pre-specified analysis cohorts (HC, $N=188$ complete cases; ALL, $N=246$ complete cases). \textbf{(b)} The eight quantitative MRI biomarkers extracted at the C2--C3 level, spanning four acquisition families: T2w cross-sectional area (T2w-CSA) from T2-weighted structural imaging; grey matter CSA (GM-CSA) from T2\textsuperscript{*}-weighted imaging; magnetisation transfer ratio (MTR) and saturation (MTsat); and fractional anisotropy (FA), mean (MD), axial (AD), and radial (RD) diffusivity from diffusion tensor imaging.}
\label{fig:baseline}
\end{figure}

\subsection{Vendor-Effect Removal}

All five methods reduced vendor $\eta^2$ by approximately three orders of magnitude, to ${\sim}10^{-4}$ (mean post-harmonisation $\eta^2$: ComBat $=4.1\times10^{-4}$, ComBat-joint $=4.9\times10^{-4}$, CovBat $=2.5\times10^{-4}$, RELIEF $=4.4\times10^{-4}$, LME/FE $=7.9\times10^{-4}$; Figure~2a,b, Table~2). For reference, the $\eta^2$ expected under the null hypothesis of no vendor effect is $(k-1)/(N-1)=2/187\approx0.011$; post-harmonisation values lie far below this chance level, so the effective reduction relative to chance expectation is 92.6--97.6\% across methods (LME/FE lowest at 92.6\%, CovBat highest at 97.6\%). Multivariate vendor structure became indistinguishable from chance (all PERMANOVA $p=1.0$ in both cohorts). These near-zero values should, however, be read as a process check rather than an independent validation: vendor is the harmonisation batch variable, and each method explicitly removes vendor location and scale, so a correctly implemented pipeline is expected to drive $\eta^2$ toward zero by construction. Machine-learnability (below) provides the stricter, non-degenerate criterion. Covariance structure was preserved differentially across methods (Supplementary Fig.~S4): CovBat and RELIEF produced near-identical cross-metric correlation matrices to each other (Frobenius distance between methods: 0.09 in HC, 0.08 in ALL), whereas ComBat, ComBat-joint, and LME/FE remained closest to the original covariance (distances to the original matrix of 1.55, 1.57, and 1.56 in HC, respectively), reflecting minimal covariance distortion by location--scale adjustment.

\begin{figure}[H]
\centering
\includegraphics[width=\textwidth]{figures/Fig2.pdf}
\caption{Harmonisation performance by method and endpoint (healthy control cohort; the baseline row shows the Original, unharmonised data). \textbf{(a)} Vendor effect size ($\eta^2$, eta-squared from one-way ANOVA) before and after harmonisation, shown on logarithmic scale. \textbf{(b)} Multivariate vendor structure (PERMANOVA $R^2$, permutational multivariate analysis of variance) on logarithmic scale. \textbf{(c)} Age-related biological signal (semi-partial $R^2$) before and after harmonisation. \textbf{(d)} Sex-related biological signal (semi-partial $R^2$) before and after harmonisation.}
\label{fig:performance}
\end{figure}

\begin{table}[H]
\centering
\scriptsize
\setlength{\tabcolsep}{2.5pt}
\caption{Harmonisation performance summary (HC cohort, complete-case $N=188$). All methods were fitted and evaluated on the same unified complete-case sample. Values are means across 8 biomarkers. $\eta^2_{\text{vendor}}$: post-harmonisation vendor ANOVA effect size; $R^2_{\text{PERMANOVA}}$: multivariate vendor structure (equals the mean per-metric $\eta^2$ by the z-score Euclidean identity); $r_{\text{pre,post}}$: subject-ranking preservation; Age $R^2$: semi-partial age variance explained. LOVO columns show leave-one-vendor-out external validation (PERMANOVA $R^2$ with z-scored Euclidean distance on the pooled within-fold sample, mean across 3 held-out-vendor folds); LOVO ratio: LOVO $\eta^2$ relative to internal $\eta^2$. Values in parentheses are 95\% percentile bootstrap CIs ($B=2{,}000$, stratified by vendor).}
\label{tab:results}
\begin{tabular}{lSSlSSS}
\toprule
Method & {$\eta^2_{\text{vendor}}$} & {$R^2_{\text{PERMANOVA}}$} & {$r_{\text{pre,post}}$ (95\% CI)} & {Age $R^2$} & {LOVO $\eta^2$} & {LOVO ratio} \\
\midrule
Original (baseline) & 3.15e-01 & 0.315 & 1.000 & 0.002 & {---} & {---} \\
ComBat & 4.09e-04 & 0.0004 & 0.789 (0.741--0.825) & 0.009 & 2.46e-03 & 6.0$\times$ \\
ComBat-joint & 4.91e-04 & 0.0005 & 0.806 (0.760--0.839) & 0.009 & 4.42e-04 & 0.9$\times$ \\
CovBat & 2.54e-04 & 0.0003 & 0.782 (0.735--0.816) & 0.010 & 1.01e-03 & 4.0$\times$ \\
RELIEF & 4.36e-04 & 0.0004 & 0.791 (0.743--0.824) & 0.009 & 1.87e-01 & 428$\times$ \\
LME/FE & 7.93e-04 & 0.0008 & 0.816 (0.771--0.848) & 0.009 & 2.01e-01 & 253$\times$ \\
\bottomrule
\end{tabular}
\end{table}

\subsection{Machine-Learnability of the Vendor Signature}

Unharmonised biomarker profiles allowed near-perfect machine detection of scanner manufacturer (macro one-vs-rest AUC $=0.986\pm0.002$ in the ALL cohort; $0.983\pm0.002$ in HC), confirming that the vendor signature dominates the multivariate feature space. After harmonisation---refitted within each cross-validation fold to prevent leakage of harmonisation parameters (Methods)---detectability decreased substantially but remained above chance for every method (permutation reference $0.495\pm0.038$): CovBat $0.629\pm0.020$, RELIEF $0.690\pm0.017$, ComBat-joint $0.704\pm0.018$, ComBat $0.706\pm0.016$, and LME/FE $0.811\pm0.015$ (ALL cohort, $N=246$; HC values were concordant: 0.670, 0.733, 0.756, 0.751, and 0.851, respectively). Expressed against the chance-adjusted detectability margin, location--scale harmonisation removed 36\% (LME/FE) to 73\% (CovBat) of the machine-detectable vendor information; the method ordering parallels correction depth, the FE fallback being a pure location shift without scale adjustment. Per-vendor detectability after ComBat ranged from 0.67 (Philips) to 0.80 (GE) (Supplementary Table~S8). Complementarily, pathology classification improved from AUC $0.520\pm0.024$ (unharmonised) to $0.582$--$0.625$ across methods (ComBat highest, $0.625\pm0.018$; point AUC 0.64, 95\% CI 0.55--0.72; within-resample paired comparison vs.\ Original: mean $\Delta$AUC $+0.105$, $p<0.001$ for all five methods by paired $t$ and Wilcoxon tests; the unpaired Hanley--McNeil test on point AUCs did not reach significance, $p=0.07$, reflecting the power gained by fold-level pairing), and sex classification remained stable ($0.646$--$0.670$), indicating that harmonisation removed vendor confounding while unmasking---rather than distorting---biological signal.

\subsection{Biological-Signal Recovery}

Before harmonisation, no biomarker showed a significant age association in either cohort (HC: all $p\geq0.34$; ALL: all $p\geq0.19$). After harmonisation, significant age associations emerged selectively in the diffusion metrics: MD and AD in both cohorts (all methods; HC $p=0.006$--$0.049$, ALL $p\leq0.0019$) and additionally RD in ALL ($p\leq0.007$); structural and magnetisation-transfer metrics remained non-significant (all $p\geq0.18$). In mean terms, age $R^2$ increased 5.6--6.2-fold in HC (from $1.5\times10^{-3}$ to $8.7\times10^{-3}$--$9.5\times10^{-3}$) and 4.8--5.0-fold in ALL (from $3.2\times10^{-3}$ to $1.5\times10^{-2}$--$1.6\times10^{-2}$); these fold changes should be read alongside the absolute values and the per-metric significance above, since a ratio of two near-zero variances is itself unstable. Sex $R^2$ showed a similar pattern (1.1--1.3-fold increase). Absolute age effects were larger in ALL, consistent with greater biological heterogeneity, but remained below 2\% of variance.

\subsection{Subject-Level Preservation and the Pareto Frontier}

Methods diverged on subject-ranking preservation (HC cohort, mean $r$ across methods: LME/FE $=0.816$, ComBat-joint $=0.806$, RELIEF $=0.791$, ComBat $=0.789$, CovBat $=0.782$). This ranking was highly concordant between cohorts (Spearman $\rho=0.90$), with a single adjacent exchange of ranks (ComBat and RELIEF swapped ranks 3--4; ALL values: LME/FE 0.801, ComBat-joint 0.787, ComBat 0.774, RELIEF 0.772, CovBat 0.763). Although the marginal bootstrap 95\% CIs for $r_{\text{pre,post}}$ overlap across methods (Table~2), a within-resample paired bootstrap ($B=2{,}000$) confirmed that the preservation hierarchy is statistically reliable: 8 of 10 method pairs were significantly separated in HC (all $p\leq0.003$) and 9 of 10 in ALL, with LME/FE significantly highest and CovBat significantly lowest (the only non-separated pairs were ComBat vs.\ CovBat and ComBat vs.\ RELIEF in HC, and ComBat vs.\ RELIEF in ALL; Supplementary Table~S7). In the HC cohort, a Pareto frontier emerged, with methods trading greater multivariate decontamination for lesser subject-level preservation (Figure~3). Because post-harmonisation PERMANOVA $R^2$ values are near-degenerate (Section~3.2), we repeated the frontier ordering using leakage-free vendor detectability (Section~3.3) as the decontamination axis: the corner solutions were preserved (CovBat: maximal removal; LME/FE: maximal preservation), and ComBat and ComBat-joint were indistinguishable on removal in HC (0.751 vs.\ 0.756, within cross-validation noise) with ComBat-joint showing higher preservation (0.806 vs.\ 0.789); in the pathology-inclusive cohort, ComBat-joint weakly dominated ComBat on both axes (lower detectability, 0.704 vs.\ 0.706, and higher preservation, 0.787 vs.\ 0.774), although the removal difference again lies within cross-validation noise. LME/FE's top rank partly reflects its minimal correction: the FE fallback applies a pure location shift without scale adjustment, leaving subject-level variance essentially untouched. Under balanced conditions (Design C, Siemens restricted to 40 subjects per seed), the preservation hierarchy reversed: ComBat-family methods converged to near-perfect preservation ($r=0.99$--$1.00$; CovBat $0.991$, others $\geq0.998$) while LME/FE remained at $r=0.843$. The main-analysis preservation gradient therefore partly reflects real-world vendor imbalance (Siemens 130/188, 69\% of HC) rather than intrinsic method properties, and the Pareto frontier should be interpreted as descriptive of this dataset's imbalance structure rather than a stable method ranking.

\begin{figure}[H]
\centering
\includegraphics[width=\textwidth]{figures/Fig3.pdf}
\caption{Pareto frontier showing the trade-off between multivariate vendor-effect removal (PERMANOVA $R^2$, permutational multivariate analysis of variance) and subject-level preservation (Pearson correlation $r_{\text{pre,post}}$ between pre- and post-harmonisation biomarker vectors). \textbf{(A)} Healthy control cohort ($N=188$); \textbf{(B)} pathology-inclusive cohort ($N=246$). A monotone frontier is present in the HC cohort; in the pathology-inclusive cohort the ordering is non-monotone (Section~3.5). Each point represents one harmonisation method. Lower right indicates maximal vendor-effect removal; upper left indicates maximal subject-level preservation.}
\label{fig:pareto}
\end{figure}

\subsection{LOVO External Validation}

ComBat and ComBat-joint generalised to unseen manufacturers without significant degradation (PERMANOVA $p\geq0.97$ and $p=1.0$ across the three folds, respectively), although ComBat's point estimate rose to LOVO $\eta^2=2.5\times10^{-3}$ (LOVO/internal ratio 6.0), reflecting estimation on the pooled within-fold sample; ComBat-joint remained at $4.4\times10^{-4}$ (ratio 0.9). CovBat showed moderate degradation ($\eta^2=1.0\times10^{-3}$, ratio 4.0). RELIEF and LME/FE failed external validation ($\eta^2_{\text{LOVO}}$: 0.19 and 0.20; PERMANOVA $p<0.01$ in all folds; ratios 428 and 253; Figure~4a--d). This pattern is explained by structural constraints: RELIEF's ICA unmixing matrix is non-transferable across vendor sets, and LME/FE cannot estimate offsets for unseen vendors. Notably, CovBat and RELIEF---the two methods that project the held-out vendor onto the training PCA/ICA basis---showed test-vendor pre/post correlations below 1 ($r=0.85$--$0.97$), whereas the affine z-score methods preserved held-out-vendor subject rankings exactly by construction.

\begin{figure}[H]
\centering
\includegraphics[width=\textwidth]{figures/Fig4.pdf}
\caption{Leave-one-vendor-out (LOVO) external validation. \textbf{(a)} Vendor effect size ($\eta^2$, eta-squared from one-way ANOVA) by method and held-out vendor, shown on logarithmic scale. \textbf{(b)} Multivariate vendor structure (PERMANOVA $R^2$, permutational multivariate analysis of variance) with significance indicators (asterisks denote *$p<0.05$, **$p<0.01$, ***$p<0.001$; ns = not significant). \textbf{(c)} Age-related biological signal (semi-partial $R^2$) recovery. \textbf{(d)} Subject-level preservation (Pearson $r_{\text{pre,post}}$) on unseen vendors. PERMANOVA $R^2$ equals the mean per-metric $\eta^2$ under the z-scored Euclidean distance (Methods), so panels (a) and (b) are numerically identical except for the significance annotation. For the unharmonised baseline (Original) no transformation is fitted, so its pooled within-fold $R^2$ ($0.315$) is identical across folds by construction.}
\label{fig:lovo}
\end{figure}

\subsection{Structural Constraints}

Three methodological observations converge on a single structural limitation of the dataset: only 3 vendors, with sites perfectly nested (43 sites in the dataset, 40 in the analysis sample). The LME random-effects covariance matrix was singular in 56\% of cases because 3 vendor levels provide insufficient information for stable REML variance-component estimation ($k\geq5$ recommended; Table~3). RELIEF's ICA found no vendor-associated components because ComBat-joint preprocessing had already removed vendor mean/scale effects, leaving insufficient residual vendor structure for ICA to detect. These constraints are inherent to the retrospective design; they do not invalidate the internal benchmark findings but bound their generalisability to similarly structured consortia.

\begin{table}[H]
\centering
\caption{LME convergence status by biomarker and cohort (complete-case unified fits). FE = fixed-effects fallback (singular random-effects covariance matrix); LME = converged linear mixed-effects model. LME converged in 7 of 16 combinations (44\%); only MD and AD achieved stable LME fits in both cohorts. All 9 FE-fallback cases produced post-harmonisation $\eta^2$ reductions exceeding 99\%, confirming that FE centring is sufficient for vendor-effect removal even when REML fails. The single converged-LME outlier (HC, T2w-CSA; post $\eta^2=4.2\times10^{-3}$) reflects an unstable variance-component estimate rather than FE fallback.}
\label{tab:lme}
\begin{tabular}{llSSSl}
\toprule
Cohort & Biomarker & {$N$} & {$\eta^2$ (pre)} & {$\eta^2$ (post)} & Status \\
\midrule
\multirow{8}{*}{HC} 
  & T2w-CSA   & 188 & 0.0530 & 4.19e-03 & \textbf{LME} \\
  & GM-CSA    & 188 & 0.1761 & 5.09e-04 & FE \\
  & MTR       & 188 & 0.2052 & 1.01e-04 & FE \\
  & MTsat     & 188 & 0.3353 & 4.78e-05 & FE \\
  & FA        & 188 & 0.2902 & 1.28e-04 & FE \\
  & MD        & 188 & 0.4646 & 3.21e-04 & \textbf{LME} \\
  & AD        & 188 & 0.7003 & 5.33e-04 & \textbf{LME} \\
  & RD        & 188 & 0.2927 & 5.07e-04 & FE \\
\midrule
\multirow{8}{*}{ALL} 
  & T2w-CSA   & 246 & 0.0441 & 3.65e-05 & FE \\
  & GM-CSA    & 246 & 0.1673 & 2.06e-05 & FE \\
  & MTR       & 246 & 0.2428 & 2.07e-05 & FE \\
  & MTsat     & 246 & 0.3596 & 1.53e-04 & \textbf{LME} \\
  & FA        & 246 & 0.3116 & 3.55e-04 & \textbf{LME} \\
  & MD        & 246 & 0.5064 & 2.88e-04 & \textbf{LME} \\
  & AD        & 246 & 0.7362 & 4.68e-04 & \textbf{LME} \\
  & RD        & 246 & 0.3199 & 4.72e-04 & FE \\
\bottomrule
\end{tabular}
\end{table}

% ===================================================================
\section{Discussion}

This study provides, to our knowledge, the first systematic benchmark of statistical harmonisation methods for multi-vendor spinal cord qMRI biomarkers. Six principal findings emerge.

First, baseline vendor effects in spinal cord qMRI are large ($\eta^2=0.31$, PERMANOVA $R^2=0.315$)---comparable to or exceeding brain imaging reports [9--11]---and dominate biological variation by orders of magnitude. For perspective, these effects approach the magnitude of pathology-related changes in quantitative metrics reported after spinal cord injury (14--20\% in R1 and magnetisation transfer saturation [7]), underscoring that unharmonised vendor offsets can masquerade as---or obscure---disease signal. Physical harmonisation, such as vendor-specific recalibration of magnetisation transfer metrics [7], offers a complementary acquisition-level route but requires prospective protocol control; statistical post-hoc harmonisation, as benchmarked here, applies retrospectively to existing consortia data. This justifies harmonisation as a default preprocessing step for multi-centre spinal cord studies---a need echoed by recent European quality initiatives [6] and by multicentre MRI harmonisation recommendations [21].

Second, all five methods reduce vendor effects to statistical negligibility on the variance-decomposition metrics each method explicitly optimises ($\eta^2$, PERMANOVA), whereas leakage-free machine-learnability---a criterion no method optimises directly---reveals substantial but incomplete removal (AUC 0.63--0.81 across methods, pathology-inclusive cohort; HC 0.67--0.85; versus 0.99 unharmonised and 0.50 chance). We note that evaluating vendor detectability after full-sample harmonisation, as is common practice, is optimistically biased because the harmonisation parameters absorb test-fold structure; our fold-internal protocol (Methods) quantifies the deployment-relevant residual. The choice among methods therefore hinges not on whether they work, but on the downstream analytic goal. For subject-level inference (individual prediction, longitudinal change detection), LME/FE or ComBat-joint may be preferred. For pooled-group analyses (case-control separation, multivariate classification), CovBat or RELIEF offer maximal multivariate purity---motivated by evidence that machine learning can exploit residual covariance differences undetectable by univariate ComBat [11], and that latent scanner structure may persist after mean/scale removal [12]. Our machine-learnability results quantify this hierarchy directly in spinal cord data: despite reducing vendor $\eta^2$ to ${\sim}10^{-4}$, LME/FE left the most detectable residual signature (AUC 0.81), whereas CovBat retained the least (0.63). The residual detectability of ComBat-harmonised profiles (AUC 0.71) is broadly consistent with the value reported for ComBat in brain imaging [11] (0.66 before covariance harmonisation), a cross-modality consistency suggesting that the residual signature left by location/scale adjustment is a general property rather than a cohort-specific artefact.

Third, the Pareto frontier between decontamination and preservation is reproducible across cohorts but is partly driven by real-world vendor imbalance: under balanced subsampling (Design C) the preservation hierarchy reversed, with ComBat-family methods converging to near-perfect preservation. The frontier should therefore be read as a descriptive characterisation of this dataset's imbalance structure rather than a stable method ranking, and the decontamination--preservation trade-off it visualises remains a useful framing for method selection.

Fourth, LOVO external validation differentiates methods sharply: ComBat and ComBat-joint generalise to unseen vendors without significant degradation, while RELIEF and LME/FE fail. Held-out-vendor degradation was markedly smaller for ComBat-joint than for per-metric ComBat (LOVO/internal $\eta^2$ ratio 0.9 vs.\ 6.0), suggesting that joint empirical Bayes estimation stabilises vendor offsets when these must be estimated from two training vendors only. For consortia planning to add new vendor scanners, ComBat-joint therefore provides a robust default choice. Out-of-sample harmonisation for new sites is an active methodological frontier: ComBat-Predict provides a principled framework for applying trained harmonisation models to unseen sites without re-accessing reference data [16], and our LOVO results independently support ComBat-family---in particular joint---harmonisation for this deployment scenario.

Fifth, the structural constraints revealed by this benchmark---vendor-site nesting, LME singular covariance in 56\% of cases, and RELIEF ICA inactivation---reflect the 3-vendor, perfectly-nested spine-generic design. These are not methodological failures but intrinsic dataset properties that bound interpretability: reported vendor effects cannot be partitioned into vendor-level versus site-level components, and LME/RELIEF's degradation in low-vendor-count settings is expected given their modelling assumptions. Larger consortia with more vendors and crossed designs would likely resolve these constraints.

Sixth, biological-signal amplification after harmonisation was selective and interpretable: age associations absent at baseline (all $p\geq0.19$) became significant in the diffusion metrics (Section~3.4; mean 4.8--6.2-fold across cohorts), and pathology classification improved (AUC 0.52 to 0.62). This demonstrates that harmonisation does not merely remove noise---it unmasks signal previously obscured by vendor variance. The concurrent pathology-classification gain mirrors the disease-classification gains reported after harmonisation in brain imaging [11], extending that observation to spinal cord biomarkers. This may have implications for clinical studies seeking to detect disease-related biomarker changes across sites.

\textbf{Limitations.} (i) Vendor-site perfect nesting precludes decomposition of vendor-level versus site-level effects; our findings reflect the composite ``vendor-in-site'' effect relevant to practical consortia. (ii) The dataset covers only three vendors with substantial imbalance (Siemens: 166/246, 67\% of the ALL cohort; 130/188, 69\% of HC); extrapolation to other manufacturers requires caution. (iii) LME converged in only 7 of 16 biomarker$\times$cohort combinations (44\%); the reported LME/FE preservation advantage ($r=0.82$) may reflect the minimal correction of the FE fallback rather than LME shrinkage. (iv) RELIEF's ICA inactivation ($q\geq0.99$ for all components) may be dataset-specific; Zhang et al.\ [12] demonstrated RELIEF's effectiveness in a three-site DTI study where latent vendor structure persisted after ComBat preprocessing---suggesting that consortia with stronger residual vendor-covariance structure may benefit from RELIEF, whereas our 43-site design with ComBat-joint preprocessing may have already exhausted this structure. (v) LOVO validation used z-score matching for unseen-vendor transformation---the standard ComBat-from-template approach [9]---but does not test prospective acquisition on genuinely new hardware; principled out-of-sample frameworks such as ComBat-Predict [16] offer a natural extension for prospective new-site harmonisation. (vi) Fold-internal machine-learnability assumes vendor labels are observed at deployment (the standard assumption of ComBat-style harmonisation); residual vendor detectability remained above chance for every method (AUC up to 0.85 for LME/FE), so harmonised pooling is not vendor-blind, and pooled analyses should continue to report vendor composition. (vii) This analysis was not pre-registered; however, all harmonisation methods were selected and evaluated on a pre-specified endpoint framework prior to analysis, and the complete analytical pipeline is publicly documented in the accompanying code repository, ensuring full computational reproducibility.

\textbf{Conclusion.} Statistical harmonisation reduces vendor effects to statistical negligibility on variance-decomposition metrics and removes the majority---though not all---of the machine-detectable vendor signature in multi-centre spinal cord qMRI. It is recommended as a default preprocessing step, with residual vendor-related confounding reported alongside pooled analyses. ComBat-joint provides the best overall balance: within-sample vendor removal was comparable to per-metric ComBat ($\eta^2=4.9$ vs.\ $4.1\times10^{-4}$, both far below the chance level) and the two methods were indistinguishable on leakage-free vendor detectability (AUC 0.704 vs.\ 0.706 in ALL), yet ComBat-joint degraded markedly less on unseen manufacturers (LOVO/internal ratio 0.9 vs.\ 6.0) and preserved subject-level information slightly better ($r=0.806$ vs.\ 0.789 in HC). ComBat-joint is therefore recommended for broad multi-centre deployment; per-metric ComBat remains appropriate when individual biomarkers are harmonised in isolation.

% ===================================================================
\section*{References}

\begin{enumerate}[nosep,leftmargin=*]
\item Cohen-Adad J, Alonso-Ortiz E, Abramovic M, et al. Open-access quantitative MRI data of the spinal cord and reproducibility across participants, sites and manufacturers. Sci Data. 2021;8:219. DOI: 10.1038/s41597-021-00941-8
\item Freund P, Boller V, Emmenegger TM, et al. Quantifying neurodegeneration of the cervical cord and brain in degenerative cervical myelopathy. Eur J Neurol. 2024;31(7):e16297. DOI: 10.1111/ene.16297
\item Valo\v{s}ek J, B\'edard S, Ke\v{r}kovsk\'y M, et al. A database of the healthy human spinal cord morphometry in the PAM50 template space. Imaging Neurosci. 2024;2:imag-2-00075. DOI: 10.1162/imag\_a\_00075
\item Cohen-Adad J, Alonso-Ortiz E, Abramovic M, et al. Generic acquisition protocol for quantitative MRI of the spinal cord. Nat Protoc. 2021;16:4611--4632. DOI: 10.1038/s41596-021-00588-0
\item Prados F, Ashburner J, Blaiotta C, et al. Spinal cord grey matter segmentation challenge. NeuroImage. 2017;152:312--329. DOI: 10.1016/j.neuroimage.2017.03.010
\item Busoni S, Pichiecchio A, Cristiano L, et al. The MR quality landscape in Europe. Insights Imaging. 2026;17(1):125. DOI: 10.1186/s13244-026-02280-x
\item Leutritz T, Seif M, Helms G, et al. Multiparameter mapping of relaxation (R1, R2*), proton density and magnetization transfer saturation at 3 T: A multicenter dual-vendor reproducibility and repeatability study. Hum Brain Mapp. 2020;41(15):4232--4247. DOI: 10.1002/hbm.25122
\item Johnson WE, Li C, Rabinovic A. Adjusting batch effects in microarray expression data using empirical Bayes methods. Biostatistics. 2007;8(1):118--127. DOI: 10.1093/biostatistics/kxj037
\item Fortin JP, Cullen N, Sheline YI, et al. Harmonization of cortical thickness measurements across scanners and sites. NeuroImage. 2018;167:104--120. DOI: 10.1016/j.neuroimage.2017.11.024
\item Fortin JP, Parker D, Tun\c{c} B, et al. Harmonization of multi-site diffusion tensor imaging data. NeuroImage. 2017;161:149--170. DOI: 10.1016/j.neuroimage.2017.08.047
\item Chen AA, Beer JC, Tustison NJ, et al. Mitigating site effects in covariance for machine learning in neuroimaging data. Hum Brain Mapp. 2022;43(4):1179--1195. DOI: 10.1002/hbm.25688
\item Zhang R, Oliver LD, Voineskos AN, et al. RELIEF: A structured multivariate approach for removal of latent inter-scanner effects. Imaging Neurosci. 2023;1:1--16. DOI: 10.1162/imag\_a\_00011
\item Bayer JMM, Thompson PM, Ching CRK, et al. Site effects how-to and when: an overview of retrospective techniques to accommodate site effects in multi-site neuroimaging analyses. Front Neurol. 2022;13:923988. DOI: 10.3389/fneur.2022.923988
\item Da-Ano R, Visvikis D, Hatt M. Harmonization strategies for multicenter radiomics investigations. Phys Med Biol. 2020;65:24TR02. DOI: 10.1088/1361-6560/aba798
\item Beer JC, Tustison NJ, Cook PA, et al. Longitudinal ComBat: a method for harmonizing longitudinal multi-scanner imaging data. NeuroImage. 2020;220:117129. DOI: 10.1016/j.neuroimage.2020.117129
\item Xin Y, Gardner M, Tustison NJ, et al. ComBat-Predict enhances generalizability of neuroimaging models to new sites. Hum Brain Mapp. 2026;47:e70546. DOI: 10.1002/hbm.70546
\item Middleton DM, Li Y, Chen A, et al. Harmonization of multi-site diffusion tensor imaging data for cervical and thoracic spinal cord at 1.5 and 3 T. Sci Rep. 2023;13:19809. DOI: 10.1038/s41598-023-46465-6
\item Jodoin P, Edde M, Girard G, et al. Challenges and best practices when using ComBat to harmonize diffusion MRI data. Sci Rep. 2025;15:41508. DOI: 10.1038/s41598-025-25400-x
\item De Leener B, Lévy S, Dupont SM, et al. SCT: Spinal Cord Toolbox, an open-source software for processing spinal cord MRI data. NeuroImage. 2017;145:24--43. DOI: 10.1016/j.neuroimage.2016.10.009
\item Anderson MJ. A new method for non-parametric multivariate analysis of variance. Austral Ecol. 2001;26:32--46. DOI: 10.1111/j.1442-9993.2001.01070.pp.x
\item De Stefano N, Battaglini M, Pareto D, et al. MAGNIMS recommendations for harmonization of MRI data in MS multicenter studies. NeuroImage Clin. 2022;34:102972. DOI: 10.1016/j.nicl.2022.102972
\end{enumerate}

\end{document}
